\newif\ifuseseminar
\useseminartrue
\input{../../latex_common/header.tex.frag}

\title{Prova de Física Matemática II -- Delta de Dirac, Funções de Green e Propagadores}

\begin{document}

\begin{enumerate}

	\item Considere a função degrau de Heaviside $\theta(x)$ e a distribuição
	      delta de Dirac $\delta(x)$.
	      \begin{enumerate}
		      \item Partindo do Teorema Fundamental do Cálculo, mostre que
		            $$
			            \frac{\dd}{\dd x}\theta(x-x_0) = \delta(x-x_0).
		            $$
		      \item Usando o resultado do item anterior, calcule
		            $\dfrac{\dd}{\dd x}\big[x\,\theta(x)\big]$ e
		            $\dfrac{\dd^2}{\dd x^2}\big[x\,\theta(x)\big]$. Interprete
		            fisicamente o resultado da segunda derivada.
		      \item Considere a gaussiana
		            $$
			            G_\sigma(x) = \frac{1}{\sqrt{2\pi\sigma^2}}\, e^{-x^2/2\sigma^2}.
		            $$
		            Mostre que $G_\sigma(x)$ é normalizada para qualquer $\sigma>0$
		            e, usando a série de Taylor de $f(x)$ em torno de $x=0$, mostre
		            que
		            $$
			            \lim_{\sigma\to0}\int_{-\infty}^{\infty} G_\sigma(x)\, f(x)\,\dd x = f(0).
		            $$
	      \end{enumerate}

	\item Considere o operador $\dfrac{\dd^2}{\dd x^2}$ atuando sobre funções
	      definidas em toda a reta real.
	      \begin{enumerate}
		      \item Verifique, por substituição direta, que
		            $$
			            G(x,x_0) = \frac{1}{2}|x-x_0|
		            $$
		            satisfaz $\dfrac{\dd^2}{\dd x^2}G(x,x_0) = \delta(x-x_0)$.
		      \item Mostre que $G(x,x_0)$ não é única: somando uma função linear
		            $a+bx$ à solução, a equação continua satisfeita. Explique,
		            sem fazer contas, que tipo de condição de contorno fixaria
		            essas constantes na prática, dando um exemplo físico.
	      \end{enumerate}

	\item Considere a função de mundo de Synge $\sigma(\vec x,\vec x_0) =
	      |\vec x-\vec x_0|^2/2$ e o laplaciano $N$-dimensional $\nabla_N^2$.
	      \begin{enumerate}
		      \item Mostre que o ansatz $G = \sigma^\lambda$ satisfaz
		            $\nabla_N^2 G = 0$ (para $\vec x\neq\vec x_0$) somente se
		            $\lambda=0$ ou $\lambda = 1-N/2$, e que a solução não trivial
		            em $N=3$ reproduz o potencial coulombiano,
		            $G\propto 1/|\vec x-\vec x_0|$.
		      \item Mostre que a solução em lei de potência falha em $N=2$
		            (justifique) e verifique, por substituição direta, que
		            $G = \ln\sigma$ satisfaz $\nabla^2 G = 0$ nesse caso.
	      \end{enumerate}

	\item Para o oscilador harmônico ($m=1$), $y''+\omega^2y=0$, o par
	      $V=(y,q)^T$, com $q=y'$, satisfaz $V'=\Omega H V$, onde
	      $\Omega=\begin{pmatrix}0&1\\-1&0\end{pmatrix}$. A solução complexa
	      normalizada (i.e., $\langle V,V\rangle\equiv iV^\dagger\Omega V=1$)
	      é dada por
	      $$
		      V(x) = \frac{1}{\sqrt{2\omega}}\begin{pmatrix}1\\-i\omega\end{pmatrix} e^{-i\omega x}.
	      $$
	      A função de Green causal (retardada) de
	      $\left(\dfrac{\dd^2}{\dd x^2}+\omega^2\right)G(x,x_0)=\delta(x-x_0)$
	      é a componente $y$ de
	      $$
		      U(x) = \Theta(x-x_0)\big[\alpha\, V(x-x_0) + \alpha^* V^*(x-x_0)\big],
	      $$
	      que satisfaz $U'-\Omega H U = \delta(x-x_0)\, U_0$, com
	      $U_0 = \alpha V(0) + \alpha^* V^*(0)$, para uma constante $\alpha$
	      apropriada.
	      \begin{enumerate}
		      \item Usando $V(0) = (1,-i\omega)^T/\sqrt{2\omega}$, escreva
		            $U_0$ explicitamente em termos de $\alpha$ e $\omega$.
		      \item O impulso instantâneo da fonte $\delta(x-x_0)$ altera
		            apenas o momento $q$, não a posição $y$, no instante
		            $x=x_0$: ou seja, $U_0$ deve ter componente $y$ nula e
		            componente $q$ igual a $1$. Use essas duas condições para
		            determinar $\alpha$.
		      \item Substituindo o $\alpha$ encontrado em $U(x)$, mostre que
		            sua componente $y$ reproduz
		            $$
			            G(x,x_0) = \Theta(x-x_0)\, \frac{\sin[\omega(x-x_0)]}{\omega},
		            $$
		            a função de Green causal do oscilador forçado.
	      \end{enumerate}

\end{enumerate}

\end{document}
